\chapter{Introdução}
\label{cap:01}

Colecionar é uma prática presente em diferentes fases da vida e nos mais diversos interesses: há quem reúna livros lidos, filmes e séries assistidos, jogos concluídos, miniaturas de carros, moedas ou discos. Mais do que acumular objetos ou títulos, a coleção envolve escolher o que entra nela, organizá-la segundo critérios próprios e, muitas vezes, exibi-la a outras pessoas. Com a digitalização do cotidiano, essas etapas passaram a acontecer também em meio digital, e o ambiente em que a coleção é mantida pode tanto favorecer quanto dificultar as práticas de aquisição, curadoria e exibição que dão sentido a ela \cite{Watkins2015}.

Registrar uma coleção não se limita a listar o que se possui ou o que se consumiu. Para muitos colecionadores, o registro funciona como memória: a data em que um filme foi assistido, a nota dada a uma temporada de série, as impressões sobre um livro ou os detalhes de uma peça adquirida permitem revisitar a coleção mais tarde, comparar opiniões antigas com as atuais e perceber a própria trajetória. A forma como essas informações são apresentadas também importa, pois a coleção é, em boa medida, uma expressão de quem a mantém.

Nesse contexto, este trabalho apresenta uma plataforma \textit{web} personalizável para criação e compartilhamento de coleções digitais. O acervo do usuário é organizado em três níveis: \textbf{Categorias}, que agrupam temas amplos (como ``Filmes'' ou ``Carros''); \textbf{Coleções}, subdivisões de uma categoria onde os registros são cadastrados e onde se definem a privacidade e o layout; e \textbf{Itens}, que representam cada unidade da coleção. A apresentação dos itens é definida por um editor de layouts no estilo de ferramentas como o Canva, no qual o usuário posiciona e redimensiona Elementos de seis tipos (texto, imagem, forma, ícone, classificação e data) e monta modelos reutilizáveis para suas coleções. Completam a proposta a personalização global da interface, por meio de paletas de cores e tipografia, e uma dimensão social, na qual o usuário pode adicionar amigos, seguir coleções de outras pessoas e acompanhar suas atualizações em um \textit{feed}.

\section{Justificativa}

A ideia deste trabalho surgiu de uma prática pessoal de registro de leituras. Ao concluir cada livro, anotavam-se detalhes como a data de término e as impressões sobre a história, com o objetivo de consultá-los depois, ao reler a obra ou simplesmente ao relembrar o que já havia sido lido. Voltar a esses registros permitia comparar a opinião de uma leitura anterior com a percepção atual sobre o enredo e os personagens.

Ao longo dos anos, diferentes meios foram experimentados para manter esse registro, sem que nenhum atendesse plenamente à necessidade. Modelos de acompanhamento de leituras encontrados na internet precisavam ser copiados à mão para o papel, o que dava trabalho e raramente alcançava o visual desejado. As planilhas eletrônicas resolviam o armazenamento, mas ofereciam poucas possibilidades de personalização, reduzindo cada registro a uma linha de tabela. Tampouco se encontrou um aplicativo dedicado que unisse o registro organizado à liberdade de definir sua aparência.

A mesma necessidade apareceu, em seguida, em outros tipos de coleção: filmes, séries, jogos e \textit{hobbies} em geral. Cada um deles pede informações diferentes (um filme pode ter diretor e data em que foi assistido; um carro de coleção, fabricante, escala e ano), e as ferramentas de campos fixos não se adaptam bem a essa variedade. Faltava, portanto, um espaço em que o colecionador pudesse decidir não apenas \textit{o que} registrar, mas \textit{como} cada registro seria apresentado, expressando sua personalidade por meio de cores, fontes e do layout das fichas.

Essa motivação encontra respaldo na literatura sobre personalização de interfaces. \citeonline{Marathe2011} observaram que, ao personalizar aspectos como fontes e cores, os usuários passam a ver a interface como reflexo de sua identidade, e que esse sentimento reforça a percepção de controle sobre a interação. Uma plataforma que une a organização de coleções à personalização visual das fichas atende, assim, tanto à necessidade prática de registrar e consultar o acervo quanto ao desejo de que esse acervo tenha a cara de quem o mantém.

\section{Objetivos}

Esta seção apresenta o objetivo geral do trabalho e os objetivos específicos que o compõem.

\subsection{Objetivo Geral}

Desenvolver uma plataforma \textit{web} que permita ao usuário criar e organizar coleções pessoais de qualquer natureza em categorias, coleções e itens, personalizar a apresentação de cada item por meio de layouts reutilizáveis e compartilhar suas coleções com outros usuários.

\subsection{Objetivos Específicos}
\begin{itemize}
	\item Levantar e especificar os requisitos funcionais e não funcionais da plataforma, bem como seus casos de uso;
	\item Elaborar o modelo de domínio do sistema e, a partir dele, o projeto de interface e o projeto de dados;
	\item Implementar a organização do acervo em categorias, coleções e itens, com controle de privacidade no nível da coleção;
	\item Construir um editor visual de layouts baseado em Elementos de seis tipos, que permita associar layouts a coleções e personalizar a ficha de itens individualmente;
	\item Oferecer a personalização global da interface por meio de paletas de cores e tipografia;
	\item Implementar as funcionalidades sociais de amizade, acompanhamento de coleções, \textit{feed} de atualizações e notificações;
	\item Verificar o funcionamento do sistema por meio de testes automatizados e testes manuais.
\end{itemize}

\section{Problema}

As ferramentas usualmente empregadas para registrar coleções pessoais, como planilhas eletrônicas, aplicativos de listas e cadernos, tratam cada registro como um conjunto de campos fixos com aparência padronizada. Essa rigidez dificulta a adaptação a acervos de naturezas distintas, que exigem informações diferentes, e não permite ao colecionador decidir como cada registro é apresentado. Em consequência, quem deseja manter coleções variadas precisa recorrer a várias ferramentas ou aceitar registros que não refletem a forma como gostaria de ver seu acervo.

O problema que este trabalho se propõe a resolver é, portanto, a ausência de uma plataforma que reúna, em um só lugar, a organização de coleções de qualquer tipo, a liberdade de definir a estrutura e a aparência de cada registro e a possibilidade de compartilhar essas coleções com outras pessoas.

\section{Organização Deste Trabalho}

Este texto está organizado em doze capítulos. O Capítulo~\ref{cap:02} reúne a fundamentação teórica que sustenta o trabalho, abordando engenharia de
requisitos, modelagem conceitual, arquitetura de aplicações \textit{web} e personalização de interfaces. O Capítulo~\ref{cap:03} descreve os materiais e os métodos empregados, da pesquisa bibliográfica à estratégia de testes, e apresenta
o equipamento e as ferramentas utilizados no desenvolvimento. O Capítulo~\ref{cap:04} traz a visão panorâmica do produto, com sua perspectiva, suas funções,
os perfis de usuário atendidos e as restrições que delimitam o escopo. O Capítulo~\ref{cap:05} especifica os requisitos do usuário e do sistema, os casos de
uso e o modelo de domínio, e é nele que se definem os quatro conceitos centrais da plataforma. O Capítulo~\ref{cap:06} trata do projeto de software, com o
projeto de interface e o projeto de dados, detalhando o mapeamento objeto-relacional e a estrutura das tabelas do banco. O Capítulo~\ref{cap:07} descreve a
implementação, incluindo o ambiente de desenvolvimento, a organização do código e as decisões tomadas durante a construção do sistema. O
Capítulo~\ref{cap:08} apresenta os testes realizados e seus resultados. O Capítulo~\ref{cap:09} trata da implantação da aplicação. O Capítulo~\ref{cap:10}
discute os resultados obtidos à luz dos objetivos estabelecidos. O Capítulo~\ref{cap:11} analisa os trabalhos e as ferramentas relacionados à proposta. Por fim, o Capítulo~\ref{cap:12} apresenta as conclusões e aponta possibilidades de
trabalhos futuros.
